% UD7 — Dispensa
% PROTOTIPO del template: per ora contiene solo le due parti che mancavano
% ai PDF UD7 (Lambert-Beer e radioattività). Le sezioni esistenti (onde EM,
% spettro, fotone, ottica) si riscrivono nel ticket "UD7 Radiazioni: Dispensa".
\dispensa

\begin{tcolorbox}[fisbox={formula}{formulaBg},title={Obiettivi dell'unità}]
Al termine di questa parte sai:
\begin{itemize}
  \item calcolare quanta radiazione attraversa uno strato di materiale con la legge di Lambert-Beer e lo spessore emivalente;
  \item descrivere un nucleo instabile e i decadimenti $\alpha$, $\beta$ e $\gamma$, scrivendo il nucleo figlio;
  \item usare la legge del decadimento radioattivo, l'attività e l'emivita.
\end{itemize}
\end{tcolorbox}

% -------------------------------------------------------------------------
\section{Assorbimento della radiazione: la legge di Lambert-Beer}
\label{sec:ud7-lambert-beer}

Quando un fascio di radiazione attraversa un materiale, una parte dei fotoni viene assorbita o deviata, e il fascio che esce è meno intenso di quello che entra. Succede alla luce in una soluzione colorata, ai raggi X nel corpo del paziente, ai raggi $\gamma$ in una parete di piombo.

\subsection{Ogni strato toglie la stessa frazione}

Immagina il materiale diviso in tanti strati sottili uguali. Ogni strato non toglie una quantità fissa di intensità, ma una \textbf{frazione fissa} di quella che gli arriva. Se il primo centimetro lascia passare metà della radiazione, il secondo centimetro lascia passare metà di quella metà, cioè un quarto del totale, e così via. Una diminuzione di questo tipo si descrive con una funzione \textbf{esponenziale}:

\begin{formula}[Legge di Lambert-Beer]
I(x) = I_0 \, e^{-\mu x}
\end{formula}

$I_0$ è l'intensità che entra, $I(x)$ quella che esce dopo uno spessore $x$, e $\mu$ è il \textbf{coefficiente di attenuazione lineare}, che si misura in \si[per-mode=power]{\metre\tothe{-1}} (o, più spesso, in \si[per-mode=power]{\centi\metre\tothe{-1}}). Il prodotto $\mu x$ è un numero puro. Il coefficiente $\mu$ dipende dal materiale e dall'energia dei fotoni: per i raggi X il piombo e l'osso hanno $\mu$ molto più grande dei tessuti molli. La legge, con un solo valore di $\mu$, vale per un fascio \textbf{monocromatico}, con fotoni tutti della stessa energia: per un fascio con energie diverse, come quello di una radiografia, è un'approssimazione.

\begin{richiamo}{esponenziali e logaritmi}
Il numero di Nepero vale $e \approx 2{,}72$. Il logaritmo naturale è l'operazione inversa dell'esponenziale: $\ln(e^{a}) = a$ e $e^{\ln b} = b$. I valori da ricordare:
\[
  \ln 2 \approx 0{,}69 \qquad \ln 10 \approx 2{,}3 \qquad e^{-1} \approx 0{,}37 \qquad \log_{10} 2 \approx 0{,}30
\]
Proprietà che servono di continuo: $e^{a}\,e^{b} = e^{a+b}$, quindi $(e^{-a})^n = e^{-na}$; e $\log_{10}(10^n) = n$. Per stimare le potenze di 2 senza calcolatrice: $2^{10} = 1024 \approx 10^3$.
\end{richiamo}

\subsection{Lo spessore emivalente}

L'esponenziale è scomodo da calcolare a mano. Conviene allora usare lo \textbf{spessore emivalente} $x_{1/2}$: lo spessore che dimezza l'intensità. Ponendo $I = I_0/2$ nella legge di Lambert-Beer si ottiene $e^{-\mu x_{1/2}} = 1/2$, cioè $\mu x_{1/2} = \ln 2$:

\begin{formula}[Spessore emivalente]
x_{1/2} = \frac{\ln 2}{\mu} \approx \frac{0{,}69}{\mu}
\qquad\qquad
I = \frac{I_0}{2^n} \quad \text{dopo } n = \frac{x}{x_{1/2}} \text{ spessori emivalenti}
\end{formula}

Con lo spessore emivalente i conti diventano potenze di 2: uno spessore emivalente lascia passare $1/2$, due lasciano passare $1/4$, tre $1/8$, dieci circa un millesimo.

\begin{figure}[h]
\centering
% Intensità trasmessa in funzione dello spessore (Lambert-Beer)
\begin{tikzpicture}
\begin{axis}[
  width=0.62\linewidth, height=5.2cm,
  axis lines=left, xmin=0, xmax=4.4, ymin=0, ymax=1.12,
  xlabel={spessore $x$}, ylabel={$I/I_0$},
  xtick={0,1,2,3,4}, xticklabels={$0$,$x_{1/2}$,$2x_{1/2}$,$3x_{1/2}$,$4x_{1/2}$},
  ytick={0,0.125,0.25,0.5,1}, yticklabels={$0$,$\tfrac18$,$\tfrac14$,$\tfrac12$,$1$},
  tick label style={font=\small}, label style={font=\small},
]
  \addplot[formula, very thick, domain=0:4.4, samples=80] {2^(-x)};
  \foreach \n/\v in {1/0.5, 2/0.25, 3/0.125} {
    \edef\tmp{\noexpand\draw[dashed, gray] (axis cs:\n,0) -- (axis cs:\n,\v) -- (axis cs:0,\v);}\tmp
  }
\end{axis}
\end{tikzpicture}

\caption{Intensità trasmessa in funzione dello spessore. A ogni spessore emivalente $x_{1/2}$ l'intensità si dimezza: la curva non arriva mai a zero.}
\label{fig:ud7-lambert-beer}
\end{figure}

\begin{esempio}
Per i raggi X di una radiografia, il tessuto molle ha uno spessore emivalente di circa \SI{3}{\centi\metre}. Quale frazione dell'intensità attraversa un addome di \SI{24}{\centi\metre}? Quanto vale $\mu$?

Lo spessore contiene $n = \SI{24}{\centi\metre} / \SI{3}{\centi\metre} = 8$ spessori emivalenti, quindi
\[
  \frac{I}{I_0} = \frac{1}{2^8} = \frac{1}{256} \approx 0{,}004 = 0{,}4\,\%
\]
Il coefficiente di attenuazione è
\[
  \mu = \frac{\ln 2}{x_{1/2}} = \frac{0{,}69}{\SI{3}{\centi\metre}} = \SI[per-mode=power]{0,23}{\centi\metre\tothe{-1}}
\]
\end{esempio}

\begin{attenzione}[l'attenuazione non è lineare]
Raddoppiare lo spessore \emph{non} dimezza la frazione trasmessa $I/I_0$: la eleva al quadrato. Se uno strato lascia passare $1/4$, due strati uguali lasciano passare $1/16$, non $1/8$. E non confondere l'attenuazione con la legge dell'inverso del quadrato, per cui l'intensità di una sorgente puntiforme diminuisce come $1/d^2$ con la distanza $d$ (la stessa potenza si distribuisce su una sfera sempre più grande): quella dipende dalla distanza dalla sorgente anche nel vuoto, l'attenuazione dipende dal materiale attraversato.
\end{attenzione}

\begin{quesitolampo}{ql:ud7-hvl}
Dopo \SI{2}{\centi\metre} di un certo materiale l'intensità di un fascio è ridotta a $1/4$ di quella iniziale. A quale frazione è ridotta dopo \SI{5}{\centi\metre}?
\opzioni{$1/10$}{$1/16$}{$1/25$}{$1/32$}{circa $1/6$}
\end{quesitolampo}

\subsection{Assorbanza e concentrazione}

In chimica e in laboratorio la stessa legge si scrive con i logaritmi in base 10. Si definisce l'\textbf{assorbanza} $A$ di un campione, e per una soluzione diluita l'assorbanza è proporzionale alla concentrazione $c$ della sostanza che assorbe e al cammino $\ell$ della luce nella cuvetta:

\begin{formula}[Assorbanza (legge di Beer)]
A = \log_{10} \frac{I_0}{I} = \varepsilon \, c \, \ell
\end{formula}

$\varepsilon$ è il \textbf{coefficiente di assorbimento molare}, caratteristico della sostanza e della lunghezza d'onda. Un'assorbanza $A = 1$ vuol dire che passa un decimo della luce, $A = 2$ un centesimo. Misurando l'assorbanza si ricava la concentrazione: è il principio dello spettrofotometro.

\begin{applicazione}{il saturimetro}
Il saturimetro da dito misura quanta parte dell'emoglobina è legata all'ossigeno. Un LED rosso (circa \SI{660}{\nano\metre}) e uno infrarosso (circa \SI{940}{\nano\metre}) illuminano il polpastrello, e un sensore misura la luce trasmessa. L'emoglobina ossigenata e quella non ossigenata assorbono in modo diverso alle due lunghezze d'onda; dal rapporto delle due assorbanze, applicando la legge di Lambert-Beer, lo strumento ricava la saturazione.
\end{applicazione}

\begin{approfondimento}{da dove viene l'esponenziale}
Uno strato sottilissimo di spessore $\mathrm{d}x$ toglie una frazione $\mu\,\mathrm{d}x$ dell'intensità che lo raggiunge: $\mathrm{d}I = -\mu I\,\mathrm{d}x$. La variazione di $I$ è proporzionale a $I$ stessa, e l'unica funzione con questa proprietà è l'esponenziale. Integrando:
\[
  \int_{I_0}^{I} \frac{\mathrm{d}I}{I} = -\mu \int_0^x \mathrm{d}x
  \quad\Longrightarrow\quad
  \ln\frac{I}{I_0} = -\mu x
  \quad\Longrightarrow\quad
  I = I_0\,e^{-\mu x}
\]
La stessa equazione, con il tempo al posto dello spessore, descrive il decadimento radioattivo della prossima sezione.
\end{approfondimento}

% -------------------------------------------------------------------------
\section{Radioattività e decadimenti radioattivi}
\label{sec:ud7-radioattivita}

\subsection{Nuclei e isotopi}

Il nucleo di un atomo contiene $Z$ protoni e $N$ neutroni. $Z$ è il \textbf{numero atomico} e decide di quale elemento si tratta; $A = Z + N$ è il \textbf{numero di massa}. Un nucleo si indica con il simbolo dell'elemento, $A$ in alto e $Z$ in basso: per esempio $\nuclide{12}{6}{C}$ è il carbonio con 6 protoni e 6 neutroni. Attenzione ai simboli: qui $A$ è il numero di massa, non l'assorbanza della sezione precedente.

Nuclei con lo stesso $Z$ e diverso $N$ sono \textbf{isotopi} dello stesso elemento: hanno le stesse proprietà chimiche ma massa diversa. Il carbonio ha, tra gli altri, gli isotopi $\nuclide{12}{6}{C}$ e $\nuclide{14}{6}{C}$.

Un nucleo è \textbf{instabile} quando ha troppi neutroni o troppi protoni rispetto ai nuclei stabili, oppure quando è troppo pesante. Prima o poi si trasforma spontaneamente in un nucleo più stabile emettendo radiazione: questo è il \textbf{decadimento radioattivo}. Gli isotopi con nucleo instabile si chiamano \textbf{isotopi radioattivi} o radioisotopi.

\subsection{I tipi di decadimento}

In ogni decadimento si conservano il numero di massa totale e la carica totale: la somma dei numeri in alto e la somma dei numeri in basso devono tornare uguali prima e dopo.

\begin{center}
\small
\begin{tabularx}{\linewidth}{@{}l X c c l@{}}
\toprule
Decadimento & Che cosa viene emesso & $\Delta A$ & $\Delta Z$ & Si ferma con \\
\midrule
$\alpha$ & un nucleo di elio $\nuclide{4}{2}{He}$ (2 protoni e 2 neutroni) & $-4$ & $-2$ & un foglio di carta, la pelle \\
$\beta^-$ & un elettrone (e un antineutrino): un neutrone diventa protone & $0$ & $+1$ & pochi mm di alluminio \\
$\beta^+$ & un positrone (e un neutrino): un protone diventa neutrone & $0$ & $-1$ & pochi mm di alluminio \\
$\gamma$ & un fotone di alta energia da un nucleo eccitato & $0$ & $0$ & cm di piombo, m di cemento \\
\bottomrule
\end{tabularx}

Il \textbf{positrone} $e^+$ è uguale all'elettrone ma con carica $+1$. Il \textbf{neutrino} $\nu$ e l'\textbf{antineutrino} $\bar\nu$ sono particelle senza carica e quasi senza massa: nel bilancio hanno $A = 0$ e $Z = 0$. Un nucleo è \textbf{eccitato} quando ha energia in eccesso, che cede emettendo un fotone $\gamma$.
\end{center}

Alcuni esempi:
\[
  \nuclide{226}{88}{Ra} \to \nuclide{222}{86}{Rn} + \nuclide{4}{2}{He}
  \qquad
  \nuclide{14}{6}{C} \to \nuclide{14}{7}{N} + e^- + \bar\nu
  \qquad
  \nuclide{18}{9}{F} \to \nuclide{18}{8}{O} + e^+ + \nu
\]
I raggi $\gamma$, come i raggi X, non hanno carica e non si fermano in un punto preciso: la loro intensità diminuisce secondo la legge di Lambert-Beer, ed è per questo che serve uno spessore di piombo tanto maggiore.

\begin{attenzione}[il decadimento $\alpha$ toglie 4 al numero di massa]
La particella $\alpha$ contiene 4 nucleoni, di cui 2 protoni: il numero di massa scende di 4 e il numero atomico di 2. Un errore frequente è togliere 2 a entrambi. Nei decadimenti $\beta$, invece, il numero di massa non cambia.
\end{attenzione}

\begin{quesitolampo}{ql:ud7-alfa}
Il nucleo $\nuclide{238}{92}{U}$ decade emettendo una particella $\alpha$. Il nucleo che si forma ha numero di massa \completa.
\end{quesitolampo}

\subsection{La legge del decadimento e l'emivita}

Non si può prevedere quando decadrà un singolo nucleo: si sa solo che ogni nucleo ha la stessa probabilità di decadere in ogni secondo, e che questa probabilità non dipende da quanto tempo è già passato. Ne segue che in ogni intervallo di tempo decade la stessa \emph{frazione} dei nuclei presenti. È la stessa situazione dell'assorbimento: il numero $N$ di nuclei non ancora decaduti diminuisce in modo esponenziale.

Quella probabilità di decadere in un secondo è la \textbf{costante di decadimento} $\lambda$, che si misura in \si[per-mode=power]{\second\tothe{-1}}. Se ci sono $N$ nuclei (qui $N$ non è più il numero di neutroni), in un secondo ne decadono in media $\lambda N$: è l'\textbf{attività} $\mathcal{A}$, il numero di decadimenti al secondo e si misura in becquerel: $\SI{1}{\becquerel} = 1$ decadimento al secondo. L'\textbf{emivita} $T_{1/2}$ è il tempo in cui metà dei nuclei decade: come per lo spessore emivalente, ponendo $N = N_0/2$ si ottiene $\lambda T_{1/2} = \ln 2$.

\begin{formula}[Decadimento radioattivo]
N(t) = N_0\, e^{-\lambda t}
\qquad\qquad
\mathcal{A} = \lambda N
\qquad\qquad
T_{1/2} = \frac{\ln 2}{\lambda}
\\
N = \frac{N_0}{2^n}, \quad \mathcal{A} = \frac{\mathcal{A}_0}{2^n} \quad \text{dopo } n = \frac{t}{T_{1/2}} \text{ emivite}
\end{formula}

L'attività è proporzionale a $N$, quindi diminuisce con la stessa emivita dei nuclei.

\begin{figure}[h]
\centering
% Nuclei non ancora decaduti in funzione del tempo
\begin{tikzpicture}
\begin{axis}[
  width=0.62\linewidth, height=5.2cm,
  axis lines=left, xmin=0, xmax=4.4, ymin=0, ymax=1.12,
  xlabel={tempo $t$}, ylabel={$N$},
  xtick={0,1,2,3,4}, xticklabels={$0$,$T_{1/2}$,$2T_{1/2}$,$3T_{1/2}$,$4T_{1/2}$},
  ytick={0,0.125,0.25,0.5,1}, yticklabels={$0$,$N_0/8$,$N_0/4$,$N_0/2$,$N_0$},
  tick label style={font=\small}, label style={font=\small},
]
  \addplot[attenzione, very thick, domain=0:4.4, samples=80] {2^(-x)};
  \foreach \n/\v in {1/0.5, 2/0.25, 3/0.125} {
    \edef\tmp{\noexpand\draw[dashed, gray] (axis cs:\n,0) -- (axis cs:\n,\v) -- (axis cs:0,\v);}\tmp
  }
\end{axis}
\end{tikzpicture}

\caption{Numero di nuclei non ancora decaduti in funzione del tempo. Dopo ogni emivita ne resta la metà.}
\label{fig:ud7-decadimento}
\end{figure}

\begin{esempio}
Per curare l'ipertiroidismo si somministra iodio-131, con emivita di \SI{8}{giorni}. Una capsula ha attività \SI{3,2}{\giga\becquerel}. Quanto vale l'attività dopo \SI{24}{giorni}? Quanti nuclei di iodio-131 contiene all'inizio?

\SI{24}{giorni} sono $n = 24/8 = 3$ emivite:
\[
  \mathcal{A} = \frac{\SI{3,2}{\giga\becquerel}}{2^3} = \SI{0,4}{\giga\becquerel} = \SI{400}{\mega\becquerel}
\]
Per il numero di nuclei serve $\lambda$, con l'emivita in secondi ($\SI{1}{giorno} = \SI{86400}{\second}$, e $8 \cdot 86\,400 = 691\,200 \approx \num{6,9e5}$); l'attività va in \si[per-mode=power]{\second\tothe{-1}}, con $\SI{1}{\giga\becquerel} = \SI{e9}{\becquerel}$:
\[
  \lambda = \frac{\ln 2}{T_{1/2}} = \frac{0{,}69}{8 \cdot \SI{86400}{\second}} \approx \frac{0{,}69}{\SI{6,9e5}{\second}} = \SI[per-mode=power]{1,0e-6}{\second\tothe{-1}}
  \qquad
  N_0 = \frac{\mathcal{A}_0}{\lambda} = \frac{\SI[per-mode=power]{3,2e9}{\second\tothe{-1}}}{\SI[per-mode=power]{1,0e-6}{\second\tothe{-1}}} = \num{3,2e15}
\]
\end{esempio}

\begin{attenzione}[dopo due emivite non è finito tutto]
Dopo un'emivita resta metà dei nuclei, ma dopo due emivite ne resta un quarto, non zero: ogni emivita dimezza quello che c'è, non toglie la stessa quantità. Per ridurre l'attività a meno di un millesimo servono circa 10 emivite.
\end{attenzione}

\begin{quesitolampo}{ql:ud7-emivita}
L'attività di un campione radioattivo scende da \SI{800}{\kilo\becquerel} a \SI{100}{\kilo\becquerel} in \SI{15}{giorni}. Quanto vale l'emivita?
\opzioni{\SI{1,9}{giorni}}{\SI{5}{giorni}}{\SI{7,5}{giorni}}{\SI{15}{giorni}}{\SI{120}{giorni}}
\end{quesitolampo}

\begin{applicazione}{medicina nucleare}
In una \textbf{scintigrafia} si somministra al paziente un radiofarmaco con tecnezio-99m, che emette fotoni $\gamma$ da \SI{140}{\kilo\electronvolt}: abbastanza penetranti da uscire dal corpo ed essere raccolti da una gamma camera. La sua emivita, circa \SI{6}{ore}, è abbastanza lunga per fare l'esame e abbastanza corta da limitare la dose al paziente.

Nella \textbf{PET} si usano emettitori $\beta^+$ come il fluoro-18 (emivita circa \SI{110}{minuti}). Il positrone emesso incontra quasi subito un elettrone: si annichilano e producono due fotoni da \SI{511}{\kilo\electronvolt} in direzioni opposte. Rivelando le coppie di fotoni lo strumento ricostruisce dove si trova il radiofarmaco.
\end{applicazione}

\begin{approfondimento}{la datazione con il carbonio-14}
Finché un organismo è vivo, scambia carbonio con l'ambiente e il rapporto fra $\nuclide{14}{6}{C}$ e $\nuclide{12}{6}{C}$ nei suoi tessuti resta quello dell'atmosfera. Dopo la morte lo scambio si ferma e il carbonio-14 decade con emivita di circa \num{5730} anni. Se in un reperto la frazione di carbonio-14 è $1/4$ di quella di un organismo vivo, sono passate due emivite: circa \num{11500} anni.
\end{approfondimento}

\riepilogoformule
